\begin{explanationbox}
	\textbf{\textcolor{CExplanation}{一句话直觉}}

	把两个根式看作点 $P(x,0)$ 到两个定点的距离；
	将定点放在 $x$ 轴两侧，折线拉直时最短。

	\medskip
	\textbf{\textcolor{CExplanation}{核心拆解}}
	\begin{description}[style=nextline, leftmargin=2.8em, labelsep=0.5em, itemsep=0.45em]
		\item[\textbf{根式就是距离：}]
		$\sqrt{(x-m)^2+a^2}$ 是 $P(x,0)$ 到
		$A(m,-a)$ 的距离；同理，
		$\sqrt{(x-n)^2+b^2}$ 是 $P$ 到 $B(n,b)$ 的距离。

		\item[\textbf{为什么把两点放在异侧：}]
		距离公式中的纵坐标差要平方，因此 $a$ 与 $-a$
		对应的距离相同。选取异侧的 $A$、$B$ 后，
		线段 $AB$ 必穿过 $x$ 轴。

		\item[\textbf{什么时候最短：}]
		$P$ 可以在整条 $x$ 轴上移动。
		当它恰好位于线段 $AB$ 与 $x$ 轴的交点时，
		折线 $A\to P\to B$ 变成直线段 $AB$，距离和最小。
	\end{description}

	\medskip
	\textbf{\textcolor{CExplanation}{算出的是什么}}

	此时纵向距离合计为 $a+b$，横向距离为 $|m-n|$，
	所以最短距离为
	\[
		|AB|=\sqrt{(m-n)^2+(a+b)^2}.
	\]

	\medskip
	\textbf{\textcolor{CExplanation}{使用场景}}
	\begin{itemize}[leftmargin=*, itemsep=0.5em, topsep=0.4em]
		\item 两个根式相加，且每个根号内都能配成
		$(x-\text{常数})^2+\text{正常数}$ 时，可尝试距离法。
		\item 若题目限制了 $x$ 的取值范围，应检查交点对应的
		$x$ 是否仍在允许范围内。
	\end{itemize}
\end{explanationbox}